Recent work on embodied dialogue shows that inserting \emph{positive friction} --- clarifying questions, assumption statements, reflective pauses --- before physical action can dramatically improve task success and user trust. Existing systems realize this behavior through hand-designed prompt controllers, which cannot directly optimize measurable outcomes such as safety violations or per-scenario success. We ask whether a reinforcement-learning agent can \emph{learn} when and which type of friction to apply. We cast the problem as a Markov Decision Process whose action space contains a single \texttt{execute} action and five friction types, whose observation is a 1162-dimensional concatenation of frozen sentence embeddings and scalar ambiguity features, and whose reward couples task success and disambiguation progress with a strong penalty on safety violations and a tunable per-turn cost. We train a flat 6-way softmax policy and a hierarchical gate-then-selector policy via Proximal Policy Optimization in a Gymnasium environment backed by a 2D world model and LLM-driven dialogue generation. On 50-episode evaluations, the flat policy attains $46.4\%$ task success with \emph{zero} safety violations, outperforming the hierarchical policy ($40.4\%$), a random baseline ($41.6\%$ with $116$ violations), and an always-execute baseline ($38.8\%$ with $152$ violations). Ablating the per-turn penalty reveals it to be the central lever of the trade-off between clarification and immediate execution. Our results suggest that learned friction policies can internalize safety signals that prompt-based controllers do not, and we discuss why a hierarchical factoring does not help at the five-action scale considered here.
